\documentclass[10pt,a4paper]{article}
\usepackage{amsmath,amssymb,amsthm,graphicx}
\usepackage{amsmath}
\usepackage{amsfonts}
\usepackage{amssymb}
\usepackage{amsthm}
\usepackage{enumerate}
\newtheorem{theorem}{Theorem}[section]
\newtheorem{conjecture}{Conjecture}[section]
\newtheorem{lemma}[theorem]{Lemma}
%\theoremstyle{definition}
\newtheorem{definition}[theorem]{Definition}
\newtheorem{corollary}[theorem]{Corollary}
\newtheorem{example}[theorem]{Example}

\newcommand{\summ}{\displaystyle\sum}
\newcommand{\prodd}{\displaystyle\prod}
\newcommand{\F}{\mathbb{F}}
\newcommand{\E}{\mathbf{E}}
\newcommand{\calF}{\mathcal{F}}
\newcommand{\Feps}{{\calF}_{\epsilon}}
\newcommand{\R}{\mathbb{R}}
\newcommand{\mon}{\mathcal{M}}
\newcommand{\supp}{\mathsf{supp}}
\newcommand{\wt}{\mathsf{wt}}
\newcommand{\poly}{\mathsf{poly}}
\newcommand{\one}{{\mathbf{1}}}
\newcommand{\qbin}[3]{\begin{bmatrix}{#1}\\ {#2}\end{bmatrix}_{#3}}
\title{Degree lower bound for polynomials computing square roots modulo primes congruent to 1 mod 4}
\author{Kiran Kedlaya, Swastik Kopparty}
\date{doi:10.4230/LIPIcs.CCC.2024.25; CC BY 4.0}
\begin{document}
\maketitle
\setcounter{section}{1}
\noindent\textit{Allowed standard prerequisites (editorial).}\begin{itemize}
\item Polynomial vanishing sets over finite fields and formal differentiation. Author Sections 1-3; original full bibliography retained
\end{itemize}\medskip
\begin{theorem}
\label{thm:1mod4}
 Let $p \equiv 1 \mod 4$. Then any polynomial that computes square roots in $\F_p$ has degree at least $\frac{p-1}{3}$.
\end{theorem}

\begin{proof}
Suppose $f(X)$ is of degree $d < \frac{p-1}{3}$ and computes square roots in $\F_p$.
%Then $f(0) = 0$, and thus $g(X) = f(X)/X$ is a polynomial of degree 
%$d-1$.
Then, since $X^{(p-1)/2} - 1$ is the vanishing polynomial of the set of nonzero perfect squares in $\F_p$, we have:
$$ f(X)^2 - X   \equiv 0  \mod (X^{(p-1)/2} - 1).$$

Let $A(X)$ be the polynomial of degree $2d-(p-1)/2$ such that
$$f(X)^2 - X =  A(X) \cdot (X^{(p-1)/2} - 1).$$
Let $B(X) = X - A(X)$.
Then we get:
\begin{equation}
\label{eq1}
f(X)^2 = A(X) \cdot X^{(p-1)/2}  + B(X),
\end{equation}
where:
\begin{itemize}
\item $\deg(f(X)) = d$.
 \item $\deg(A(X)), \deg(B(X)) \leq 2d - (p-1)/2$.
 \item $A(X) \neq 0$. If $A(X) = 0$ then $f(X)^2 = X$, which is impossible for a polynomial $f(X)$. 
 \item $B(X) \neq 0$. Otherwise $A(X) = X$, and $f(X)^2 = X^{(p+1)/2}$, which is possible only if $p \equiv 3 \mod 4$.
\end{itemize}

These conditions together will give us our lower bound on $d$.

Taking derivatives\footnote{Throughout this paper, we work with formal derivatives of polynomials.} of both sides of~\eqref{eq1}, we get:
\begin{align}
\label{eq2}
2f(X)f'(X) = A'(X) \cdot X^{(p-1)/2} - \frac{1}{2}A(X) X^{(p-3)/2} + B'(X).
\end{align}

Computing $2 f(X)^2 f'(X)$ in two ways using~\eqref{eq1} and~\eqref{eq2}, we get:
$$ 2 f'(X) A(X) X^{(p-1)/2} + 2 f'(X) B(X) = f(X) \left(X \cdot A'(X) - \frac{1}{2}A(X)\right) X^{(p-3)/2} + f(X) B'(X).$$
Now, using our assumption on $d$, the degrees of $2 f'(X)B(X)$ and $f(X) B'(X)$ are both at most
$3d - (p-1)/2 - 1 < (p-3)/2$, and thus taking the above equation mod $X^{(p-3)/2}$,
$$ 2f'(X) B(X) = f(X) B'(X).$$

Since $B(X) \neq 0$, we get $\frac{2 f'(X)}{f(X)} = \frac{B'(X)}{B(X)}$. Since $2\deg(f), \deg(B) < p$, by a basic property of logarithmic derivatives, this implies $f(X)^2 = \lambda B(X)$ for some nonzero $\lambda$, contradicting the fact that $A(X) \neq 0$. (See Remark 1 in Section~\ref{sec:rem} for a precise statement and a proof.)


Thus our assumption that $d < \frac{p-1}{3}$ is wrong, and the theorem follows.
\end{proof}

Let $p$ be an odd prime, and let $\F_p$ be the finite field with $p$ elements.

We say $f(X) \in \F_p[X]$ {\em computes square roots} in $\F_p$, if
for all nonzero perfect squares $a \in \F_p$, we have:
$$f(a)^2 = a.$$ In other words, for each nonzero perfect square $a \in \F_p$, $f(a)$ is one of its two square roots.

\subsection*{Logarithmic derivatives}\label{sec:rem}
 The key fact about logarithmic derivatives that we are using is that if 
 $f(X), B(X) \in \F_p[X]$, $t \cdot \deg(f), \deg(B) < p$, and:
 $$ t \frac{f'(X)}{f(X)} = \frac{B'(X)}{B(X)},$$
 then $f(X)^t = \lambda \cdot B(X)$ for some constant $\lambda \in \F_p$. 
 
 We recap a quick proof. The hypothesis implies that $\left( \frac{f(X)^t}{B(X)} \right)' = 0$,
  and thus $\frac{f(X)^t}{B(X)}$ must be a rational function in $X^p$. To see the last deduction, note that $\left(f(X)^t \cdot B(X)^{p-1} \right)' = \left( \frac{f(X)^t}{B(X)} \cdot B(X)^p \right)' = \left( \frac{f(X)^t}{B(X)} \right)' \cdot B(X)^p + \frac{f(X)^t}{B(X)} \cdot p \cdot B(X)^{p-1} \cdot B'(X) = 0 + 0 = 0$. Thus the polynomial $f(X)^t \cdot B(X)^{p-1}$ is a polynomial in $X^p$, which implies that $\frac{f(X)^t}{B(X)} = 
 \frac{f(X)^t \cdot B(X)^{p-1}}{B(X)^p} =
 \frac{f(X)^t \cdot B(X)^{p-1}}{B(X^p)}$ is a rational function in $X^p$. 
 Once we know that $\frac{f(X)^t}{B(X)}$ is a rational function in $X^p$, our assumption about the degrees implies the result. 

\begin{thebibliography}{VZGG13}

\bibitem[ADN03]{ADN-sqrt}
Simon~Joseph Agou, Marc Del{\'e}glise, and Jean-Louis Nicolas.
\newblock Short polynomial representations for square roots modulo p.
\newblock {\em Designs, Codes and Cryptography}, 28(1):33--44, 2003.

\bibitem[AS97]{ASlowdegree}
Sanjeev Arora and Madhu Sudan.
\newblock Improved low-degree testing and its applications.
\newblock In {\em Proceedings of the twenty-ninth annual ACM symposium on
  Theory of computing}, pages 485--495, 1997.

\bibitem[Bir00]{Biro}
Andr{\'a}s Bir{\'o}.
\newblock On polynomials over prime fields taking only two values on the
  multiplicative group.
\newblock {\em Finite Fields and Their Applications}, 6(4):302--308, 2000.

\bibitem[BS96]{BachShallit}
Eric Bach and Jeffrey~Outlaw Shallit.
\newblock {\em Algorithmic number theory: Efficient algorithms}, volume~1.
\newblock MIT press, 1996.

\bibitem[Car11]{Carella-sqrt}
NA~Carella.
\newblock Formulas for the square root modulo p.
\newblock {\em arXiv preprint arXiv:1101.4605}, 2011.

\bibitem[CKL15]{CKL-nthroot}
Seunghwan Chang, Bihtnara Kim, and Hyang-Sook Lee.
\newblock Polynomial representations for n-th roots in finite fields.
\newblock {\em Journal of the Korean Mathematical Society}, 52(1):209--224,
  2015.

\bibitem[CS00]{CS-dlog}
Don Coppersmith and Igor Shparlinski.
\newblock On polynomial approximation of the discrete logarithm and the
  diffie—hellman mapping.
\newblock {\em Journal of Cryptology}, 13:339--360, 2000.

\bibitem[GR05]{GR-listrecovery}
Venkatesan Guruswami and Atri Rudra.
\newblock Limits to list decoding reed-solomon codes.
\newblock In {\em Proceedings of the thirty-seventh annual ACM symposium on
  Theory of computing}, pages 602--609, 2005.

\bibitem[GS98]{RSlist}
Venkatesan Guruswami and Madhu Sudan.
\newblock Improved decoding of reed-solomon and algebraic-geometric codes.
\newblock In {\em Proceedings 39th Annual Symposium on Foundations of Computer
  Science (Cat. No. 98CB36280)}, pages 28--37. IEEE, 1998.

\bibitem[Kal91]{KaltofenHIT}
Erich Kaltofen.
\newblock Effective noether irreducibility forms and applications.
\newblock In {\em Proceedings of the twenty-third annual ACM symposium on
  Theory of Computing}, pages 54--63, 1991.

\bibitem[Mas84]{abcthm2}
Richard~Clive Mason.
\newblock {\em Diophantine equations over function fields}, volume~96.
\newblock Cambridge University Press, 1984.

\bibitem[MW86]{MW-dlog}
Gary Mullen and David White.
\newblock A polynomial representation for logarithms in gf (q).
\newblock {\em Acta arithmetica}, 3(47):255--261, 1986.

\bibitem[Sch06]{Schmidtbook}
Wolfgang~M Schmidt.
\newblock {\em Equations over finite fields: an elementary approach}, volume
  536.
\newblock Springer, 2006.

\bibitem[Sha73]{ShanksTonelli1}
Daniel Shanks.
\newblock Five number-theoretic algorithms.
\newblock In {\em Proceedings of the Second Manitoba Conference on Numerical
  Mathematics (Winnipeg), 1973}, 1973.

\bibitem[Sho09]{Shoup}
Victor Shoup.
\newblock {\em A computational introduction to number theory and algebra}.
\newblock Cambridge university press, 2009.

\bibitem[Sto81]{abcthm1}
W~Wilson Stothers.
\newblock Polynomial identities and hauptmoduln.
\newblock {\em The Quarterly Journal of Mathematics}, 32(3):349--370, 1981.

\bibitem[Ton91]{ShanksTonelli2}
Alberto Tonelli.
\newblock Bemerkung {\"u}ber die aufl{\"o}sung quadratischer congruenzen.
\newblock {\em Nachrichten von der K{\"o}nigl. Gesellschaft der Wissenschaften
  und der Georg-Augusts-Universit{\"a}t zu G{\"o}ttingen}, 1891:344--346, 1891.

\bibitem[VZGG13]{GerhardGathen}
Joachim Von Zur~Gathen and J{\"u}rgen Gerhard.
\newblock {\em Modern computer algebra}.
\newblock Cambridge university press, 2013.

\bibitem[WB83]{BerlekampWelch}
Lloyd~R Welch and Elwyn~R Berlekamp.
\newblock Error correction for algebraic block codes.
\newblock {\em US patent}, (4,633,470), 1983.

\bibitem[Win02]{W-dlog}
Arne Winterhof.
\newblock Polynomial interpolation of the discrete logarithm.
\newblock {\em Designs, Codes and Cryptography}, 25(1):63--72, 2002.

\end{thebibliography}
\end{document}
